% ---
% id: statistics/binomial-box/rem:guard-activation
% kind: rem
% title: "Regimes in which the guard acts"
% summary: >-
%   Four exact comparisons of the allocation with the binomial distribution
%   function decide whether the guard acts and whether it costs more than one
%   count. It acts even at ordinary budgets, on any outcome whose positive
%   expected count is at most its tail budget.
% prominence: standard
% status: accepted
% uses: [prop:rule-exact, lem:guard, statistics/cor:tail-forms, statistics/lem:exact-strict-tails, statistics/thm:binomial-tails, def:interval, def:endpoint-rule, def:construction, statistics/def:matched-alphabet]
% mentions: [thm:monotonicity-limits, rem:guard, req:rule-in-force, req:request-defaults, protocol-compiler/def:spec, protocol-compiler/ass:coin]
% assumes: [A-TSM, A-ALPHABET]
% refs: []
% provenance: newly-derived
% lean:
% ---

\begin{remark}\label{rem:guard-activation}
This remark reads the guard on the exact rule of
Definition~\ref{def:interval}, written for a single honest value
$p_{\hat c}=\honest(\hat c)$, the degenerate envelope; for a general envelope
every statement holds side by side, by the side-wise reading of
Definition~\ref{def:endpoint-rule}.

The four exact conditions \eqref{eq:guard-lo0}--\eqref{eq:guard-up1} of
Proposition~\ref{prop:rule-exact}\,(vi) say where the guard acts and what it
costs there. By \eqref{eq:guard-lo0} and \eqref{eq:guard-up0} the guard is
entirely inert as soon as each allocated tail budget is strictly below the
mass the count law places at or below $\lfloor np_{\hat c}\rfloor$,
respectively at or above $\lceil np_{\hat c}\rceil$; by \eqref{eq:guard-lo1}
and \eqref{eq:guard-up1} it costs at most one count under the correspondingly
weaker conditions. All four are comparisons of an allocated budget against a
value of the binomial distribution function, hence decidable exactly in the
closed forms \eqref{eq:tails-used} of Corollary~\ref{cor:tail-forms}: each
is the negation of a comparison of Lemma~\ref{lem:exact-strict-tails}\,(ii)
at the count next to the one it names, or a boundary value of
Theorem~\ref{thm:binomial-tails} where that count leaves $\{0,\dots,n\}$,
with no machinery beyond what
Definition~\ref{def:interval} already requires. The two
pairs fail in different regimes, however, and are not to be read together.

\emph{A cost of more than one count} --- failure of \eqref{eq:guard-lo1} or of
\eqref{eq:guard-up1} --- requires a single outcome's single tail to carry a
budget at least the mass the count law places up to one step beyond
$np_{\hat c}$, that is
$\varepsilon^{-}(\hat c)\ge\Pr[\Ncount{\hat c}\le\lfloor np_{\hat c}\rfloor+1]$
respectively
$\varepsilon^{+}(\hat c)\ge\Pr[\Ncount{\hat c}\ge\lceil np_{\hat c}\rceil-1]$.
At the ends of the range that cannot happen at all: if $np_{\hat c}<1$ then
$\lceil np_{\hat c}\rceil-1\le0$ and
$\Pr[\Ncount{\hat c}\ge\lceil np_{\hat c}\rceil-1]=1>\varepsilon^{+}(\hat c)$,
and if $\lfloor np_{\hat c}\rfloor\ge n-1$ then
$\lfloor np_{\hat c}\rfloor+1\ge n$ and
$\Pr[\Ncount{\hat c}\le\lfloor np_{\hat c}\rfloor+1]=1>\varepsilon^{-}(\hat
c)$, both using $\varepsilon^{\pm}(\hat c)\le\ecomp<1$. In the example of
Remark~\ref{rem:guard} the budget is indeed near-unit: $\varepsilon^{-}(\hat
c)=\tfrac9{10}$ against $\Pr[\Ncount{\hat c}\le6]=\tfrac{53}{64}$. That is
also the regime of the $\ecomp=\tfrac45$-on-one-outcome counterexample of
Theorem~\ref{thm:monotonicity-limits}\,(v), and under an allocation that concentrates a
near-unit budget on one or two constrained outcomes the enlargement is to be
read off \eqref{eq:guard-cost}, not estimated by a rule of thumb.

\emph{Activation of the guard at all} --- failure of \eqref{eq:guard-lo0} or
of \eqref{eq:guard-up0} --- is a different matter: it is not confined to
near-unit budgets. It occurs at ordinary budgets and under the engine's
default \tref[uniform rule]{uniform allocation rule}, for any outcome whose expected count over
the whole run is positive and at or below the tail budget allocated to it.
Precisely, if $0<np_{\hat c}\le\varepsilon^{+}(\hat c)$ --- which forces
$np_{\hat c}<1$ and hence $\lceil np_{\hat c}\rceil=1$, since
$\varepsilon^{+}(\hat c)<1$ --- then Bernoulli's inequality
$(1-p_{\hat c})^{n}\ge1-np_{\hat c}$ gives
\begin{equation*}
  \Pr\big[\Ncount{\hat c}\ge\lceil np_{\hat c}\rceil\big]
  =1-(1-p_{\hat c})^{n}\ \le\ np_{\hat c}\ \le\ \varepsilon^{+}(\hat c),
\end{equation*}
so \eqref{eq:guard-up0} fails, $\kappa^{+}(\hat c)=0$, and the guard raises
the emitted endpoint to $k^{+}(\hat c)=1$. Symmetrically, if
$0<n(1-p_{\hat c})\le\varepsilon^{-}(\hat c)$ then $\lfloor
np_{\hat c}\rfloor=n-1$ and $\Pr[\Ncount{\hat c}\le n-1]=1-p_{\hat c}^{\,n}\le
n(1-p_{\hat c})\le\varepsilon^{-}(\hat c)$ --- the same inequality applied at
$1-p_{\hat c}$ --- so \eqref{eq:guard-lo0} fails, $\kappa^{-}(\hat c)=n$, and
the guard lowers the emitted endpoint to $k^{-}(\hat c)=n-1$. A concrete
instance is an \tref{admissible input} with the engine's default rule and
default allocation (Requirements~\ref{req:rule-in-force}
and~\ref{req:request-defaults}): $\Chatset=\{\hat c_0,\hat c_1\}$, $n=10$,
$\honest=(\tfrac{999}{1000},\tfrac1{1000})$, $\Chatcon=\Chatset$,
$\ecomp=\tfrac1{10}$ and the uniform allocation
$\varepsilon^{\pm}(\hat c)=\ecomp/(2|\Chatcon|)=\tfrac1{40}$; there
$np_{\hat c_1}=n(1-p_{\hat c_0})=\tfrac1{100}<\tfrac1{40}$, so the guard is
active on the upper side of $\hat c_1$ and on the lower side of $\hat c_0$
--- at a budget nowhere near $1$, and with the budget spread uniformly rather
than concentrated. Nor is this a small-$n$ artefact: at $n=10^{6}$ with
$\ecomp=\tfrac1{10}$ and $|\Chatcon|=5$ the uniform rule allocates
$\varepsilon^{\pm}(\hat c)=10^{-2}$, and a constrained outcome of honest
probability $10^{-9}$ has $np_{\hat c}=10^{-3}\le\varepsilon^{+}(\hat c)$, so
the guard forces $q^{+}(\hat c)=1/n$ where the raw quantile interval would
have given $0$. Under A-TSM (\xref{reductions/ass:tsm}), this
module's alphabet being that of the consuming declaration
(Definition~\ref{def:matched-alphabet}), the acceptance alphabet is the
joint outcome, carrying every per-round public choice: one made at a
measuring station as a component of the outcome by A-TSM, and one on the
letter side as a component of the letter
(\xref{protocol-compiler/def:spec}\,(a)), drawn as declared by A-COIN
(\xref{protocol-compiler/ass:coin}); so outcomes of
expected count at or below one are a routine feature of a fine-grained
$\Chatset$ rather than an exotic one. What the guard costs on them is
nevertheless bounded, by both measures: exactly one count --- the relevant
condition of the previous paragraph holds in each case,
$\lceil np_{\hat c}\rceil-1\le0$ giving \eqref{eq:guard-up1} on the upper
side and $\lfloor np_{\hat c}\rfloor+1\ge n$ giving \eqref{eq:guard-lo1} on
the lower side --- and at most $1/n$ from the honest value by
\eqref{eq:guard-from-p}.
\end{remark}
